\section{Examples using \pmt}\label{pure tableau examples section}

\par In this section we present several examples of derivations in $\pmt$. We will assume the soundness (Theorem~\ref{Soundness of PMT - theorem}) and completeness (Theorem~\ref{completeness of flt theorem}) of our system $\pmt$ for now, but this will be proven in subsequent sections. If the formula is not originally in $\Lp$, we rewrite it so that it is $\Lp$ --- Proposition~\ref{pure modal expresses basic proposition}. Generally, when we input a formula into a tableau, we first compose any composable modalities so that it is in $\LpMax$ as per Definition~\ref{maximally composed formulas definition - pure}. This is not required, and we will do an example where we don't compose all modalities.

\par The first example we will do is an equivalent form of the K-axiom in $\Lp$. We would expect the tableau to show that the K-axiom is valid. When we input the negation of a valid formula into a tableau, assuming soundness, the tableau will close.

\begin{example}
    We start by rewriting the K-axiom as a formula in $\Lp$ and then composing any modalities.
     %Ignored for First submission
%\nnote{I have used red brackets (and other colours) to make it more readable. If this is distracting, I can remove it.}
\begin{align*}
    [\textbf{1}]\tcr{(} p \to q\tcr{)} \to \tcb{(}[\textbf{1}]p \to [\textbf{1}]q\tcb{)}
    &\equiv [\textbf{1}]\tcr{(}\neg p \vee q\tcr{)} \to \tcb{(}\neg[\textbf{1}]p \vee [\textbf{1}]q\tcb{)}\\
    &\equiv \neg[\textbf{1}]\tcr{(}\neg p \vee q\tcr{)} \vee \tcb{(}\neg[\textbf{1}]p \vee [\textbf{1}]q\tcb{)}\\
    &\equiv \neg[\textbf{1}]\tcr{(}[\iota_2](\neg p, q)\tcr{)} \vee [\iota_2]\tcb{(}\neg[\textbf{1}]p, [\textbf{1}]q\tcb{)}\\
    &\equiv \neg[\textbf{1}(\iota_2)]\tcr{(}\neg p, q\tcr{)} \vee [\iota_2]\tcb{(}[\iota_1](\neg[\textbf{1}]p), [\textbf{1}]q\tcb{)}\\
    &\equiv \neg[\textbf{1}(\iota_2)]\tcr{(}\neg p, q\tcr{)} \vee [\iota_2(\iota_1,\tb{1})]\tcb{(}\neg[\textbf{1}]p, q\tcb{)}\\
    &\equiv [\iota_2]\tcg{(}\neg[\textbf{1}(\iota_2)]\tcr{(}\neg p, q\tcr{)} , [\iota_2(\iota_1,\tb{1})]\tcb{(}\neg[\textbf{1}]p, q\tcb{)}\tcg{)}\\
    &\equiv [\iota_2]\tcg{(}[\iota_1](\neg[\textbf{1}(\iota_2)]\tcr{(}\neg p, q\tcr{)}), [\iota_2(\iota_1,\tb{1})]\tcb{(}\neg[\textbf{1}]p, q\tcb{)}\tcg{)}\\
    &\equiv [\iota_2(\iota_1,\iota_2(\iota_1,\tb{1}))]\tcg{(}\neg[\textbf{1}(\iota_2)]\tcr{(}\neg p, q\tcr{)} , \neg[\textbf{1}]p, q\tcg{)}\\
\end{align*}

Now, let's input the negation of this and prove that it is valid:



\begin{center}
\renewcommand*\linenumberstyle[1]{(#1)}
\begin{tableau}
    {vertical/.style={
before drawing tree={not ignore edge, edge=draw},
},}
    [\neg\lsb\iota_2(\iota_1\coms\iota_2(\iota_1\coms\tb{1}))\rsb\tcg{(}\neg\lsb\textbf{1}(\iota_2)\rsb\tcr{(}\neg p\coms q\tcr{)} \coms \neg\lsb\textbf{1}\rsb p\coms q\tcg{)}:x
    [R_{\iota_2(\iota_1\coms\iota_2(\iota_1\coms\tb{1}))}xy_1y_2y_3 \coms \neg\neg\lsb\textbf{1}(\iota_2)\rsb\tcr{(}\neg p\coms q\tcr{)}:y_1 \coms \neg\neg\lsb\tb{1}\rsb p:y_2 \coms \neg q:y_3 ,just={$(\langle\alpha\rangle E)$ from (1)}
    [\lsb \textbf{1}(\iota_2)\rsb (\neg p\coms  q) :y_1 ,just={$(\neg E)$ from (2)}
    [\lsb\textbf{1}\rsb(\lsb\iota_2\rsb(\neg p\coms  q)) :y_1 ,just={$([\alpha(\beta)]D)$ from (3)}
    [\lsb\textbf{1}\rsb p:y_2 ,just={$(\neg E)$ from (2)}
    [R_{\iota_2}xz_1z_2 \coms \riO z_1y_1 \coms R_{\iota_2(\iota_1,\tb{1})}z_2y_2y_3 ,just={$(RR)$ from (2)}
    [R_{\iota_2}z_2w_1w_2 \coms R_{\iota_1}w_1y_2 \coms R_{\tb{1}}w_2y_3 ,just={$(RR)$ from (6)}
    [\lsb\textbf{1}\rsb p:w_1 ,just={$(\riotn x)$ from (7) and (5)}
    [\lsb\textbf{1}\rsb p:w_2 ,just={$(\riotn x)$ from (7) and (8)}
    [p:y_3 ,just={$([\alpha]E)$ from (7) and (9)}
    [\lsb\textbf{1}\rsb(\lsb\iota_2\rsb(\neg p\coms  q)) :z_1 ,just={$(\riotn x)$ from (6) and (4)}
    [\lsb\textbf{1}\rsb(\lsb\iota_2\rsb(\neg p\coms  q)) :z_2 ,just={$(\riotn x)$ from (6) and (11)}
    [\lsb\textbf{1}\rsb(\lsb\iota_2\rsb(\neg p\coms  q)) :w_2 ,just={$(\riotn x)$ from (7) and (12)}
    [\lsb\iota_2\rsb(\neg p\coms  q) :y_3 ,just={$([\alpha]E)$ from (7) and (13)}
    [\riT y_3y_3y_3 ,just={$(\riT I)$}
    [\neg p:y_3,close][q:y_3 ,close ,just={$([\alpha]E)$ from (15) and (14)}
    ]]]]]]]]]]]]]]]]
\end{tableau}
\end{center}
Here (16a) closes w.r.t.\ (10) and (16b) closes w.r.t.\ (2). Therefore, the tableau is closed, i.e.\ it could not find a counter model for the K-axiom. Hence, we can conclude that the K-axiom is valid (using soundness, Theorem~\ref{Soundness of PMT - theorem}).
\end{example}


\begin{example}
    In this example, we rewrite $\Diamond\top \to (\Box p \to \Diamond p)$ as a formula in $\LpMax$, and then produce a $\pmt$ tableau deriving it. 
    %Ignored for First submission
%\nnote{correct parsing?}

    \begin{align*}
        \neg\tcb{(}\langle\textbf{1}\rangle[\iota_0] \rightarrow \tcg{(}[\textbf{1}] p \rightarrow \langle\textbf{1}\rangle (p)\tcg{)}\tcb{)}
        &\equiv \neg\tcb{(}\neg\langle\textbf{1}\rangle[\iota_0] \vee \tcg{(}\neg[\textbf{1}] p \vee \langle\textbf{1}\rangle p\tcg{)}\tcb{)}\\
        &\equiv \neg\tcb{(}\neg\neg[\textbf{1}]\neg[\iota_0] \vee \tcg{(}\neg[\textbf{1}] p \vee \neg[\textbf{1}]\neg p\tcg{)}\tcb{)}\\
        &\equiv \neg\tcb{(}[\textbf{1}]\neg[\iota_0] \vee \tcg{(}\neg[\textbf{1}] p \vee \neg[\textbf{1}]\neg p\tcg{)}\tcb{)}\\
        &\equiv \neg\tcb{(}[\textbf{1}]\neg[\iota_0] \vee\neg[\textbf{1}] p \vee \neg[\textbf{1}]\neg p\tcb{)}\\
        &\equiv \neg\tcb{(}[\iota_3]\tcg{(}[\textbf{1}]\neg[\iota_0] ,\neg[\textbf{1}] p , \neg[\textbf{1}]\neg p\tcg{)}\tcb{)}\\
        &\equiv \neg\tcb{(}[\iota_3]\tcg{(}[\textbf{1}]\neg[\iota_0] ,[\iota_1](\neg[\textbf{1}] p) , [\iota_1](\neg[\textbf{1}]\neg p)\tcg{)}\tcb{)}\\
        &\equiv \neg\tcb{(}[\iota_3(\textbf{1},\iota_1,\iota_1)]\tcg{(}\neg[\iota_0],\neg[\textbf{1}] p , \neg[\textbf{1}]\neg p\tcg{)}\tcb{)}\\
        &\equiv \neg[\iota_3(\textbf{1},\iota_1,\iota_1)](\neg[\iota_0],\neg[\textbf{1}] p , \neg[\textbf{1}]\neg p)
    \end{align*}

    \par Although the formula has now been written to be in $\LpMax$, $\pmt$ does not handle $\iota_n$ for any $n>2$, so we need to reduce these with $\iota_1$'s and $\iota_2$'s.

    \begin{align*}
        \neg[\iota_3(\textbf{1},\iota_1,\iota_1)](\neg[\iota_0],\neg[\textbf{1}] p , \neg[\textbf{1}]\neg p)
        &\equiv \neg[\iota_2(\iota_1(\textbf{1}),\iota_2(\iota_1,\iota_1))](\neg[\iota_0],\neg[\textbf{1}] p , \neg[\textbf{1}](\neg p))\\
        &\equiv \neg[\iota_2(\textbf{1},\iota_2))](\neg[\iota_0],\neg[\textbf{1}] p , \neg[\textbf{1}]\neg p)
    \end{align*}

\par The last step is convenient, but not required.

\begin{center}
\renewcommand*\linenumberstyle[1]{(#1)}
\begin{tableau}
    {vertical/.style={
before drawing tree={not ignore edge, edge=draw},
},}
    [\neg\lsb \iota_2(\textbf{1}\coms \iota_2)\rsb (\neg\lsb \iota_0\rsb \coms  \neg\lsb \textbf{1}\rsb  p \coms \neg\lsb \textbf{1}\rsb \neg p) :x
    [R_{\iota_2(\textbf{1},\iota_2)}xy_1y_2y_3 \coms  \neg\neg\lsb \iota_0\rsb :y_1 \coms  \neg\neg\lsb \textbf{1}\rsb  p :y_2 \coms  \neg\neg\lsb \textbf{1}\rsb \neg p:y_3, just={$(\langle\alpha\rangle E)$ from (1)}
    [\lsb \textbf{1}\rsb  p :y_2 ,just={$(\neg E)$ from (2)}
    [\lsb \textbf{1}\rsb \neg p:y_3 ,just={$(\neg E)$ from (2)}
    [R_{\iota_2}xz_1z_2 \coms  R_{\textbf{1}}z_1y_1 \coms  R_{\iota_2}z_2y_2y_3 ,just={$(RR)$ from (2)}
    [R_\textbf{1}xy_1 ,just={$(R_{\iota_n}R_\alpha)$ from (5)}
    [\lsb \textbf{1}\rsb  p :z_2 ,just={$(R_{\iota_n}x)$ from (5) and (3)}
    [\lsb \textbf{1}\rsb \neg p:z_2 ,just={$(R_{\iota_n}x)$ from (5) and (4)}
    [\lsb \textbf{1}\rsb  p :x ,just={$(R_{\iota_n}x)$ from (5) and (7)}
    [\lsb \textbf{1}\rsb \neg p:x ,just={$(R_{\iota_n}x)$ from (5) and (8)}
    [ \neg p :y_1,just={$([\alpha]E)$ from (6) and (10)}
    [p:y_1, close, just={$([\alpha]E)$ from (6) and (9)}]]]]]]]]]]]]
\end{tableau}
\end{center}
\end{example}

\par Let's test the validity of a formula that does not belong in the basic modal language.

\begin{example}
    For this example, we prove the validity of $[\textbf{2}](p,q) \vee \langle\textbf{2}\rangle(\neg p, \neg q)$, but we won't be composing every possible modality. 
    %Ignored for First submission
%\nnote{Phrasing?}
    \begin{align*}
        \neg([\textbf{2}](p,q) \vee \langle\textbf{2}\rangle(\neg p, \neg q)) & \equiv \neg([\iota_2]\tcb{(}[\textbf{2}](p,q) , \langle\textbf{2}\rangle(\neg p, \neg q)\tcb{)})\\
        & \equiv \neg([\iota_2]\tcb{(}[\textbf{2}](p,q) , \neg[\textbf{2}](\neg\neg p, \neg\neg q)\tcb{)})\\
        & \equiv \neg[\iota_2]\tcb{(}[\textbf{2}](p,q) , \neg[\textbf{2}](p, q)\tcb{)}\\
    \end{align*}

\begin{center}
\renewcommand*\linenumberstyle[1]{(#1)}
\begin{tableau}
{vertical/.style={
before drawing tree={not ignore edge, edge=draw},
},}
    [\neg\lsb\iota_2\rsb(\lsb\textbf{2}\rsb(p\coms q) \coms \neg\lsb\textbf{2}\rsb(p\coms q)):x
    [R_{\iota_2}xy_1y_2 \coms \neg\lsb\textbf{2}\rsb(p\coms q):y_1 \coms \neg\neg\lsb\textbf{2}\rsb(p\coms q) :y_2 ,just={$(\langle\alpha\rangle E)$ from (1)}
    [\lsb\textbf{2}\rsb(p\coms q) :y_2 ,just={$(\neg E)$ from (2)}
    [\lsb\textbf{2}\rsb(p\coms q) :x ,just={$(R_{\iota_n}x)$ from (2) and (3)}
    [\neg\lsb\textbf{2}\rsb(p\coms q) :x ,close,just={$(R_{\iota_n}x)$ from (2)}
    ]]]]]
\end{tableau}
\end{center}

% \begin{center}
% \renewcommand*\linenumberstyle[1]{(#1)}
% \begin{tableau}
% {vertical/.style={
% before drawing tree={not ignore edge, edge=draw},
% },}
%     [\neg\lsb\iota_2\rsb(\lsb\textbf{2}\rsb(p\coms q) \coms \neg\lsb\textbf{2}\rsb(p\coms q)):x
%     [R_{\iota_2(\textbf{2}\coms\iota_1)}xy_1y_2y_3 \coms \neg p:y_1 \coms \neg q:y_2 \coms \neg\neg\lsb\textbf{2}(p\coms q)\rsb :y_3 ,just={$(\langle\alpha\rangle E)$ from (1)}
%     [\lsb\textbf{2}\rsb(p\coms q) :y_3 ,just={$(\neg E)$ from (2)}
%     [R_{\iota_2}xz_1z_2 \coms R_\textbf{2}z_1y_1y_2 \coms R_{\iota_1}z_2y_3 ,just={$(RR)$ from (2)}
%     [R_\textbf{2}xy_1y_2 ,just={$(R_{\iota_n}R_\alpha)$ from (4)}
%     [\lsb\textbf{2}\rsb(p\coms q) :z_2 ,just={$(R_{\iota_n}x)$ from (4) and (3)}
%     [\lsb\textbf{2}\rsb(p\coms q) :x ,just={$(R_{\iota_n}x)$ from (4) and (6)}
%     [p:y_1,close][q:y_2,close, just={$([\alpha]E)$ from (5) and (7)}]
%     ]]]]]]]    
% \end{tableau}
% \end{center}
    \par The node (5) closes with respect to (4). Since the tableau could not find a counter-model for $[\textbf{2}](p,q) \vee \langle\textbf{2}\rangle(\neg p, \neg q)$, then, with soundness, it is valid. Note that $[\textbf{2}](p,q) \vee \langle\textbf{2}\rangle(\neg p, \neg q)$ can be rewritten as $[\textbf{2}](p,q) \vee \neg [\textbf{2}](p, q)$, showing that $[\textbf{2}](p,q) \vee \langle\textbf{2}\rangle(\neg p, \neg q)$ follows the same scheme as the propositional formula $p \vee \neg p$. This is not a pure modal language instance of $p \vee \neg p$ since $\vee$ is not in $\Lp$.
    \par We can use this example to show that the tableau for a formula is not unique. Since there is no priority of rules that we need to use, we could have implemented the $(\langle\alpha\rangle E)$ rule to $\neg[\textbf{2}](p,q):y_1$, and the tableau would have taken more steps to close.
\end{example}

\begin{remark}
    It is appropriate here to mention that the closure of tableau is independent of the order of rule application.
\end{remark}

\par Before we go into the last example of this section, we will point out a more systematic way of keeping track of what formulas and relational atoms are in our set $\Delta$. Since we don't copy all the previous conclusions over (when displaying tableaux), a more systematic way of doing this is necessary --- especially when the tableau gets long.
\par For every tableau rule, there are labelled formulas and relational atoms that get replaced, and others get copied. When doing a tableau (using Fitting's style \cite{fitting1972tableau} to display these tableaux), we can tick off any labelled formulas/relational atoms that were replaced by the tableau rule. I.e., anything copied by the tableau rule is not ticked off. Let's see how this works when applying one rule.
\par If we have the following at a node: 
\begin{center}
\renewcommand*\linenumberstyle[1]{}
\begin{tableau}
    {vertical/.style={
before drawing tree={not ignore edge, edge=draw},
},}
    [\Delta \coms \neg\lsb \alpha\rsb (\phi_1\text{
    ,\,} \phi_2\coms  \ldots\coms  \phi_{m}):x]
\end{tableau}
\end{center}

\par Then applying the $(\langle\alpha\rangle E)$ rule, $\neg\lsb \alpha\rsb (\phi_1\text{,\,} \phi_2\coms  \ldots\coms  \phi_{m})$ is replaced by $R_\alpha xy_1\ldots y_m \coms  \neg\phi_1 :y_1 \coms  \ldots\coms \\ \neg\phi_m :y_m$ and so we can tick off $\neg\lsb \alpha\rsb (\phi_1\text{,\,} \phi_2\coms  \ldots\coms  \phi_{m})$. And in the tableau, it will look as follows:

\begin{center}
\renewcommand*\linenumberstyle[1]{}
\begin{tableau}
    {}
    [\Delta \coms \neg\lsb \alpha\rsb (\phi_1\text{
    ,\,} \phi_2\coms  \ldots\coms  \phi_{m}):x\checkmark
    [\Delta \coms R_\alpha xy_1\ldots y_m \coms  \neg\phi_1 :y_1 \coms  \ldots \coms  \neg\phi_m :y_m ,just={$(\langle\alpha\rangle E)$}
    ]]
\end{tableau}
\end{center}

\par Note: we have been doing this implicitly in the previous examples, so this is not describing a new way of doing these tableaux. So, let's do an example with this in mind:

\begin{example}\label{Open ft tableau}
With this example we will input the negation of $\neg(\langle \textbf{1} \rangle (p) \wedge (\langle\textbf{1}\rangle (q) \wedge [\textbf{1}](\neg(p \wedge q))))$ to deduce if this formula is valid. Again, we assume the soundness and completeness of $\pmt$.
\begin{align*}
    \neg\neg(\langle \textbf{1} \rangle (p) \wedge \tcr{(}\langle\textbf{1}\rangle (q) \wedge [\textbf{1}]\tcb{(}\neg(p \wedge q)\tcb{)}\tcr{)})
    &\equiv \langle \textbf{1} \rangle (p) \wedge \tcr{(}\langle\textbf{1}\rangle (q) \wedge [\textbf{1}]\tcb{(}\neg(p \wedge q)\tcb{)}\tcr{)}\\
    &\equiv \langle \textbf{1} \rangle (p) \wedge \tcr{(}\langle\textbf{1}\rangle (q) \wedge [\textbf{1}]\tcb{(}\neg p \vee \neg q\tcb{)}\tcr{)}\\
    &\equiv \langle \textbf{1} \rangle (p) \wedge \tcr{(}\langle\textbf{1}\rangle (q) \wedge [\textbf{1}]\tcb{(}[\iota_2](\neg p \coms \neg q)\tcb{)}\tcr{)}\\
    &\equiv \langle \textbf{1} \rangle (p) \wedge \tcr{(}\langle\textbf{1}\rangle (q) \wedge [\textbf{1}(\iota_2)]\tcb{(}\neg p, \neg q\tcb{)}\tcr{)}\\
    &\equiv \langle \textbf{1} \rangle (p) \wedge \langle \iota_2 \rangle\tcr{(}\langle\textbf{1}\rangle (q), [\textbf{1}(\iota_2)]\tcb{(}\neg p, \neg q\tcb{)}\tcr{)}\\
    &\equiv \langle \iota_2 \rangle\tcg{(}\langle \textbf{1} \rangle (p) , \langle \iota_2 \rangle\tcr{(}\langle\textbf{1}\rangle (q), [\textbf{1}(\iota_2)]\tcb{(}\neg p, \neg q\tcb{)}\tcr{)}\tcg{)}\\
    &\equiv \langle \iota_2 \rangle\tcg{(}\langle \textbf{1} \rangle (p) , \langle \iota_2 \rangle\tcr{(}\neg[\textbf{1}] (\neg q), [\textbf{1}(\iota_2)]\tcb{(}\neg p, \neg q\tcb{)}\tcr{)}\tcg{)}\\
    &\equiv \langle \iota_2 \rangle\tcg{(}\langle \textbf{1} \rangle (p) , \neg[\iota_2]\tcr{(}\neg\neg [\textbf{1}] (\neg q), \neg[\textbf{1}(\iota_2)]\tcb{(}\neg p, \neg q\tcb{)}\tcr{)}\tcg{)}\\
    &\equiv \langle \iota_2 \rangle\tcg{(}\neg[\textbf{1}] (\neg p) , \neg[\iota_2]\tcr{(} [\textbf{1}] (\neg q), \neg[\textbf{1}(\iota_2)]\tcb{(}\neg p, \neg q\tcb{)}\tcr{)}\tcg{)}\\
    &\equiv \neg[\iota_2] \tcg{(}\neg\neg[\textbf{1}] (\neg p) , \neg\neg[\iota_2]\tcr{(} [\textbf{1}] (\neg q), \neg[\textbf{1}(\iota_2)]\tcb{(}\neg p, \neg q\tcb{)}\tcr{)}\tcg{)}\\
    &\equiv \neg[\iota_2] \tcg{(}[\textbf{1}] (\neg p) , [\iota_2]\tcr{(}[\textbf{1}] (\neg q), \neg[\textbf{1}(\iota_2)]\tcb{(}\neg p, \neg q\tcb{)}\tcr{)}\tcg{)}\\
    &\equiv \neg[\iota_2] \tcg{(}[\textbf{1}] (\neg p) , [\iota_2]\tcr{(}[\textbf{1}] (\neg q), [\iota_1](\neg[\textbf{1}(\iota_2)]\tcb{(}\neg p, \neg q\tcb{)})\tcr{)}\tcg{)}\\
    &\equiv \neg[\iota_2] \tcg{(}[\textbf{1}] (\neg p) , [\iota_2(\textbf{1}, \iota_1)]\tcr{(}\neg q, \neg[\textbf{1}(\iota_2)]\tcb{(}\neg p, \neg q\tcb{)}\tcr{)}\tcg{)}\\
    &\equiv \neg[\iota_2(\textbf{1}, \iota_2(\textbf{1}, \iota_1))] \tcg{(}\neg p , \neg q, \neg[\textbf{1}(\iota_2)]\tcb{(}\neg p, \neg q\tcb{)}\tcg{)}\\
\end{align*}

\begin{center}
\renewcommand*\linenumberstyle[1]{(#1)}
\resizebox{160mm}{!}{
\begin{tableau}
    {vertical/.style={
before drawing tree={not ignore edge, edge=draw},
},}
    [\neg\lsb \iot{2}(\tb{1}\coms\iot{2}(\tb{1}\coms\iot{1}))\rsb(\neg p \coms \neg q \coms \neg\lsb\tb{1}(\iot{2})\rsb(\neg p\coms  \neg q)) :x \checkmark
    [R_{\iot{2} (\textbf{1},\iot{2}(\textbf{1},\iot{1}))}xy_1y_2y_3 \checkmark \coms  \neg\neg p:y_1 \checkmark  \coms  \neg\neg q:y_2 \checkmark  \coms  \neg\neg\lsb\textbf{1}(\iot{2})\rsb(\neg p \coms  \neg q):y_3 \checkmark , just={($\langle\alpha\rangle E$) from (1)}
    [p:y_1 , just={($\neg E)$ from (2)}
    [q:y_2 , just={($\neg E)$ from (2)}
    [ \lsb  1(\iot{2}) \rsb (\neg p \coms  \neg q):y_3  \checkmark , just={($\neg E)$ from (2)}
    [R_{\iot{2}}xz_1z_2 \coms R_\textbf{1}z_1y_1 \coms R_{\iot{2}(\textbf{1},\iot{1})}z_2y_2y_3 \checkmark, just={$(RR)$ from (2)}
    [R_{\iot{2}}z_2w_1w_2 \coms R_\textbf{1}w_1y_2 \coms R_{\iot{1}}w_2y_3 , just={$(RR)$ from (6)}
    [\lsb\textbf{1}\rsb (\lsb \iot{2}\rsb (\neg p \coms  \neg q) ) :y_3, just={ ([$\alpha(\beta)]D)$ from (5)}
    [\lsb\textbf{1}\rsb (\lsb \iot{2}\rsb (\neg p \coms  \neg q) ) :w_2, just={ ($R_{\iota_n}x$) from (7) and (8)}
    [\lsb\textbf{1}\rsb  (\lsb \iot{2}\rsb  (\neg p \coms  \neg q) ) :z_2, just={ ($R_{\iota_n}x$) from (7) and (9)}
    [R_\textbf{1}z_2y_2 , just={$(R_{\iota_n}R_{\alpha}$) from (7)}
    [\lsb \iot{2}\rsb  (\neg p \coms  \neg q) :y_2 , just={($[\alpha]E$) from (11) and (10)}
    [\riT y_2y_2y_2 ,just={$(\riT I)$}
    [\neg q: y_2,close][\neg p:y_2, just={$(\lsb 
    \alpha\rsb E$), (14a) closes w.r.t.\ (4)}
    [\lsb\textbf{1}\rsb  (\lsb \iot{2}\rsb  (\neg p \coms  \neg q) ) :x, just={ ($R_{\iota_n}x$) from (6) and (10)}
    [R_\tb{1}xy_1, just={$(\riotn R_\alpha)$ from (6)}
    [\lsb \iot{2}\rsb (\neg p\coms\neg q):y_1, just={$([\alpha]E)$ from (16) and (15)}
    [\riT y_1y_1y_1, just={$(\riT I)$}
    [\neg p:y_1,close][\neg q:y_1, just={$(\lsb\alpha\rsb E$), (19a) closes w.r.t.\ (3)}
    ]]]]]]]]]]]]]]]]]]]
\end{tableau}}
\end{center}

\par We continue the tableau from (19).

\begin{center}
\renewcommand*\linenumberstyle[1]{(#1)}
\begin{tableau}
    {vertical/.style={
before drawing tree={not ignore edge, edge=draw},
},line no shift=18}
    [\neg q:y_1
    [\lsb\textbf{1}\rsb (\lsb \iot{2}\rsb (\neg p \coms  \neg q) ) :w_1, just={$(\riotn x)$ from (7) and (9)}
    [\lsb\textbf{1}\rsb (\lsb \iot{2}\rsb (\neg p \coms  \neg q) ) :z_1, just={$(\riotn x)$ from (6) and (10)}
    [R_\tb{1}z_2y_1, just={$(\riotn R_\alpha)$ from (6)}
    [R_\tb{1}w_1y_1, just={$(\riotn R_\alpha)$ from (7) and (22)}
    [R_\tb{1}w_2y_1, just={$(\riotn R_\alpha)$ from (7) and (22)}
    [R_\tb{1}y_3y_1, just={$(\riotn R_\alpha)$ from (7) and (24)}
    [R_\textbf{1}w_2y_2 , just={$(R_{\iota_n}R_{\alpha}$) from (7)}
    [R_\textbf{1}y_3y_2 , just={$(R_{\iota_n}R_{\alpha}$) from (7) and (26)}
    [R_\textbf{1}z_1y_2 , just={$(R_{\iota_n}R_{\alpha}$) from (6) and (11)}
    [R_\textbf{1}xy_2 , just={$(R_{\iota_n}R_{\alpha}$) from (6) and (11)}
    ]]]]]]]]]]]
\end{tableau}
\end{center}

\par We will stop the tableau here. This is for presentation reasons. The tableau will reach all its end nodes without closing. To see this, notice that none of $(\neg E), (\langle\alpha\rangle E), (RR), ([\alpha(\beta)]D)$ are applicable to any non-ticked tableau terms occurring in the tableau. Also, for any application of $(\riotn x)$, it would repeat what is already on the branch. This is the same with the $([\alpha]E)$ rule. We can still apply the $(\riO I)$ and $(\riT I)$ rules, but this adds no `useful' information. Now the $(\riotn R_\alpha)$ rule can still be applied to the relational atoms that have $\iota_1$ or $\iot{2}$ as their underlying relation, but displaying this will result in a tableau that has more than two hundred steps. This is because $x,y_3,z_1,z_2,w_1,w_2$ are related to each other with some $\iota$-relation and so $(R_{\iotn}R_\alpha)$ will have to be applied multiple times so that every combination of $\riT ***$ and $\riO **$ can be listed --- $6\times 6\times 6>200$. Once this is done, every rule will be used up to saturation. Since this tableau is open, we can find a model to satisfy the original formula. But to say this with confidence, we need completeness.
\par We reiterate what the $\checkmark$'s indicate with an example. To create node (6), we apply (RR) using $R_{\iota_2 (\textbf{1},\iota_2(\textbf{1},\iota_1))}xy_1y_2y_3$ from the displayed node (2). This relational atom is replaced with three new relational atoms, and so $R_{\iota_2 (\textbf{1},\iota_2(\textbf{1},\iota_1))}xy_1y_2y_3$ can be ticked off. What this means is that $R_{\iota_2 (\textbf{1},\iota_2(\textbf{1},\iota_1))}xy_1y_2y_3$ is not part of the implicit $\Delta$ in node (6).
\end{example}

\par These are all the examples that we will do in this section. We have assumed soundness and completeness of $\pmt$, but the next two sections are dedicated to proving the soundness and completeness of $\pmt$.